\section{Desarrollo de las sesiones}

\begin{sesionbox}{1}{Inicio}{presentamos el reto}
\campo{Inicio (15 min).} Plano y fotografía de los paneles del gimnasio. Lectura del encargo de la dirección. Relato breve de cómo Eratóstenes midió la Tierra con una sombra. Lluvia de ideas: ¿qué no podemos medir directamente?

\campo{Desarrollo (30 min).} Cuestionario diagnóstico en GeoGebra Classroom \parencite{geogebra}: razones en el triángulo rectángulo, Pitágoras y coordenadas. Presentación del centro socio de eTwinning y de los cuatro roles de equipo.

\campo{Cierre (10 min).} ¿Qué necesitaríamos saber para responder? Cada estudiante estima a ojo la altura del edificio principal y hasta dónde llega su sombra. Preguntas y estimaciones se anotan en un mural que se revisa en S20.

\campo{Agrupamiento:} gran grupo e individual. \campo{Recursos:} proyector, portátiles. \campo{Evaluación:} diagnóstico, sin calificación; sirve para formar equipos heterogéneos y preparar el refuerzo.
\end{sesionbox}

\begin{sesionbox}{2}{Taller}{exploramos los vectores}
\campo{Inicio (15 min).} Se publican los equipos. Formación en los roles: tarjetas del turno~1, normas del equipo, fórmula de la crítica razonada y creación del tablero Kanban con la primera tarea.

\campo{Desarrollo (30 min).} Ficha F1: representar vectores, sumarlos y multiplicarlos por un número arrastrando sus extremos; observar qué ocurre con el producto escalar al cambiar el ángulo y cuándo vale cero. Un vector hace de rayo de sol y otro de normal al panel.

\campo{Cierre (10 min).} Puesta en común: ¿cuándo recibe el panel el sol de frente? Las conjeturas se formalizan en S3 y S4.

\campo{Agrupamiento:} individual con ayuda del equipo. \campo{Recursos:} GeoGebra, F1. \campo{Evaluación:} F1 (7.1, 2.1); observación (9.3).
\end{sesionbox}

\begin{sesionbox}{3}{Magistral}{el conjunto de los vectores}
\campo{Exposición (25 min).} Estructura y propiedades del conjunto de vectores. Suma y producto por un número, a partir de lo observado en F1. Módulo y coordenadas respecto a una base.

\campo{Ejercicios (20 min).} Práctica individual graduada, con corrección comentada.

\campo{Pieza del reto (10 min).} Sobre el plano a escala del centro, cada equipo escribe el vector que va del pie del edificio a una zona candidata y calcula la distancia real.

\campo{Agrupamiento:} gran grupo, individual y equipos. \campo{Recursos:} pizarra, GeoGebra en proyector, hoja de ejercicios, plano del centro. \campo{Evaluación:} seguimiento de los ejercicios, sin calificación.
\end{sesionbox}

\begin{sesionbox}{4}{Magistral}{producto escalar}
\campo{Exposición (25 min).} Definición de producto escalar, $\vec u\cdot\vec v=|\vec u||\vec v|\cos\alpha$, propiedades y expresión en coordenadas. Ángulo entre vectores, proyección ortogonal y perpendicularidad.

\campo{Ejercicios (15 min).} Práctica individual graduada, con corrección comentada.

\campo{Pieza del reto (10 min).} Ejemplo en equipo: ángulo entre el rayo solar del mediodía y la normal de un panel inclinado.

\campo{Cierre (5 min).} Cuestionario Q1 de autocorrección.

\campo{Agrupamiento:} gran grupo, individual y equipos. \campo{Recursos:} pizarra, GeoGebra en proyector, hoja de ejercicios. \campo{Evaluación:} Q1 (7.1, 5.2).
\end{sesionbox}

\begin{sesionbox}{5}{Taller}{la circunferencia goniométrica}
\campo{Inicio (5 min).} Del triángulo rectángulo a un ángulo cualquiera: ¿qué le pasa a la sombra de una varilla cuando sube el sol?

\campo{Desarrollo (40 min).} Ficha F2: construir la circunferencia goniométrica con un deslizador para el ángulo, observar seno, coseno y tangente en los cuatro cuadrantes y formular conjeturas ($\sin^2\alpha+\cos^2\alpha=1$, ángulo doble).

\campo{Cierre (10 min).} Se contrastan las conjeturas; las erróneas se analizan como oportunidad para aprender. Para casa: construir el clinómetro.

\campo{Agrupamiento:} individual y parejas. \campo{Recursos:} GeoGebra, F2. \campo{Evaluación:} F2 (3.2, 3.1); observación (3.2, 9.1).
\end{sesionbox}

\begin{sesionbox}{6}{Magistral}{razones e identidades trigonométricas}
\campo{Exposición (30 min).} Razones de un ángulo cualquiera en grados y radianes. Demostración de la identidad fundamental y de las razones del ángulo suma y diferencia.

\campo{Ejercicios (15 min).} Demostraciones guiadas y cálculo de razones con las fórmulas.

\campo{Pieza del reto (10 min).} Con $h=90^\circ-\varphi+\delta$, cada equipo calcula la altura del sol y la sombra de una varilla de 1~m en varias fechas.

\campo{Agrupamiento:} gran grupo, individual y equipos. \campo{Recursos:} pizarra, hoja de ejercicios, calculadora. \campo{Evaluación:} demostraciones escritas (3.1).
\end{sesionbox}

\begin{sesionbox}{7}{Magistral}{ángulo doble y ecuaciones trigonométricas}
\campo{Exposición (25 min).} Razones del ángulo doble y mitad. Método para resolver ecuaciones trigonométricas sencillas: soluciones en $[0^\circ,360^\circ)$ y solución general. Ejemplo resuelto: $\sin 2x=\cos x$.

\campo{Ejercicios (12 min).} Ecuaciones de dificultad creciente.

\campo{Pieza del reto (10 min).} Resolver $\sin(\beta_p+h)=1$ para hallar la inclinación del panel con incidencia perpendicular en distintas fechas.

\campo{Cierre (8 min).} Cuestionario Q2 (5 min) y relevo de roles (3 min).

\campo{Agrupamiento:} gran grupo, individual y equipos. \campo{Recursos:} pizarra, hoja de ejercicios. \campo{Evaluación:} Q2 (1.2).
\end{sesionbox}

\begin{sesionbox}{8}{Magistral}{teoremas del seno y del coseno}
\campo{Exposición (30 min).} Demostración del teorema del seno y del teorema del coseno. Criterio para elegir cada uno según los datos.

\campo{Ejercicios (15 min).} Resolución de triángulos, justificando qué teorema conviene en cada caso.

\campo{Pieza del reto (10 min).} Primera medida real: con su clinómetro y una cinta, cada equipo mide la altura del aula desde un punto de medida y la compara con el valor real. El rol de medición dirige la calibración del clinómetro.

\campo{Agrupamiento:} gran grupo, individual y equipos. \campo{Recursos:} pizarra, hoja de ejercicios, clinómetro, cinta. \campo{Evaluación:} seguimiento de los ejercicios, sin calificación.
\end{sesionbox}

\begin{sesionbox}{9}{Magistral}{caso ambiguo y altura inaccesible}
\campo{Exposición (20 min).} Caso ambiguo: dos lados y el ángulo opuesto a uno de ellos. Ejemplo resuelto: altura inaccesible desde dos puntos de medida.

\campo{Ficha (20 min).} Ficha F3: comprobar en GeoGebra, arrastrando vértices, cuándo hay dos soluciones, una o ninguna, y justificar cada caso.

\campo{Pieza del reto (10 min).} Cada equipo planifica sobre el plano los dos puntos de medida para medir el bloque de viviendas, al que no se puede llegar.

\campo{Cierre (5 min).} Cuestionario Q3.

\campo{Agrupamiento:} gran grupo, parejas y equipos. \campo{Recursos:} pizarra, GeoGebra, F3, plano del centro. \campo{Evaluación:} F3 (1.1, 1.2, 2.1); Q3 (2.1).
\end{sesionbox}

\begin{sesionbox}{10}{Magistral}{ecuaciones de la recta}
\campo{Exposición (25 min).} Ecuaciones de la recta a partir de un punto y un vector director: vectorial, paramétricas, continua, general y explícita. Vector normal. Posición relativa de puntos y rectas.

\campo{Ejercicios (20 min).} Elegir la ecuación más útil en cada caso y estudiar posiciones relativas.

\campo{Pieza del reto (10 min).} El rayo solar como recta: cada equipo escribe la del edificio elegido y decide si varios puntos del plano quedan en sombra.

\campo{Agrupamiento:} gran grupo, individual y equipos. \campo{Recursos:} pizarra, GeoGebra en proyector, plano del centro. \campo{Evaluación:} seguimiento de los ejercicios, sin calificación.
\end{sesionbox}

\begin{sesionbox}{11}{Magistral}{ángulos y distancias}
\campo{Exposición (25 min).} Incidencia, paralelismo y ortogonalidad. Ángulo entre rectas y distancia de un punto a una recta mediante el producto escalar.

\campo{Ejercicios (12 min).} Problemas métricos de dificultad creciente.

\campo{Pieza del reto (10 min).} Las filas de paneles como rectas paralelas: cada equipo compara el coeficiente del mediodía con el del IDAE y explica por qué el pliego exige más separación.

\campo{Cierre (8 min).} Cuestionario Q4 (5 min) y relevo de roles (3 min).

\campo{Agrupamiento:} gran grupo, individual y equipos. \campo{Recursos:} pizarra, GeoGebra en proyector, hoja de ejercicios. \campo{Evaluación:} Q4 (5.2).
\end{sesionbox}

\begin{sesionbox}{12}{Taller}{rayo y sombra}
\campo{Inicio (5 min).} Cada estudiante revisa la tarjeta de su nuevo rol y se recuerda qué pide el encargo.

\campo{Desarrollo (40 min).} Ficha F4 en equipo. Todos los miembros trabajan juntos; el equipo puede pedir tarjetas de ayuda, y quien termina antes aborda la ampliación.

\campo{Cierre (10 min).} Puesta en común sobre qué ecuación conviene (5 min). Cada equipo completa el croquis de su método a partir de la pieza de S9 (5 min).

\campo{Agrupamiento:} equipos. \campo{Recursos:} GeoGebra, F4, tarjetas de ayuda. \campo{Evaluación:} F4 (7.2); observación (9.1).
\end{sesionbox}

\begin{sesionbox}{13}{Taller}{refuerzo, mapa conceptual y método}
\campo{Inicio (5 min).} Devolución de los cuestionarios y reparto por criterios.

\campo{Refuerzo y ampliación (20 min).} Trabajo individual o por parejas. Quien lo necesita revisa, con ejercicios resueltos paso a paso, el criterio que indiquen sus cuestionarios: 7.1 y 5.2 (Q1), 1.2 (Q2), 2.1 (Q3) o 5.2 (Q4). El resto aborda la ampliación: panel inclinado óptimo o más de un obstáculo en el modelo.

\campo{Mapa y método (25 min).} Cada equipo elabora el mapa conceptual que relaciona vectores, trigonometría y rectas, a partir de una lista de conceptos. Mientras, el profesor pasa por cada equipo, unos 3 minutos, para aprobar su método de medida.

\campo{Cierre (5 min).} Revisión del Kanban y preparación de la salida.

\campo{Agrupamiento:} individual, parejas y equipos. \campo{Recursos:} ejercicios resueltos, plantilla del mapa, croquis. \campo{Evaluación:} mapa (5.1); segundo intento de los cuestionarios antes de S18.
\end{sesionbox}

\begin{sesionbox}{14}{Tarea}{medimos el instituto}
\campo{Preparación.} La sesión se sitúa en la franja que incluye el mediodía solar (hacia las 13:25 en la segunda semana de marzo; la hora exacta se consulta en la calculadora de la NOAA \parencite{noaa}).

\campo{Desarrollo (45 min).} En el patio: medir con clinómetro y cinta, repitiendo cada medida. Se mide un objeto accesible desde uno y desde dos puntos de medida, para comparar los métodos, y el bloque de viviendas desde dos puntos de medida. Medir la sombra de una varilla vertical de 1~m al mediodía solar y subir el dato al TwinSpace.

\campo{Si hay nubes.} Las alturas se miden igual, porque el clinómetro no necesita sol. La sombra de la varilla la mide un representante de cada equipo el primer día despejado de esa semana, y se anota la fecha. Si no hay ninguno, se usa la altura del sol que da la NOAA y se indica en el libro.

\campo{Cierre (10 min).} Cada equipo valida sus datos y detecta medidas que debe repetir.

\campo{Agrupamiento:} equipos con roles. \campo{Recursos:} clinómetro, cinta, varilla, hoja de datos. \campo{Evaluación:} observación (9.1, 9.3).
\end{sesionbox}

\begin{sesionbox}{15}{Tarea}{modelo de sombras y análisis del error}
\campo{Inicio (5 min).} Objetivos del día en el Kanban.

\campo{Desarrollo (40 min).} A partir de las piezas del reto ya resueltas: modelo a escala en GeoGebra con deslizador para la fecha, resolución de triángulos y comprobación a mano. Efecto de un grado de error y elección del resultado. Sombras en solsticios y equinoccios.

\campo{Cierre (10 min).} Prueba de reutilización: cambiar fecha y altura y comprobar que el modelo responde. Actualización del libro y relevo de roles (3 min).

\campo{Agrupamiento:} equipos. \campo{Recursos:} GeoGebra, piezas del reto, hoja de datos. \campo{Evaluación:} libro (1.1, 2.2, 7.2); observación (9.3).
\end{sesionbox}

\begin{sesionbox}{16}{Tarea}{propuesta de ampliación y coevaluación}
\campo{Inicio (5 min).} Objetivos del día en el Kanban.

\campo{Desarrollo (30 min).} Cálculo de filas con el criterio del IDAE, plano de la propuesta, número de paneles y producción con PVGIS, con la orientación del departamento de Tecnología. Comparación con los datos del socio y reflexión sobre ODS y límites del modelo.

\campo{Cierre (20 min).} El profesor modela una crítica razonada (3 min). Coevaluación del borrador entre equipos con R1 (12 min) y anotación de las críticas en el registro de cambios (5 min).

\campo{Agrupamiento:} equipos. \campo{Recursos:} GeoGebra, PVGIS, TwinSpace, R1. \campo{Evaluación:} libro y propuesta (6.1, 6.2, 7.2); coevaluación (9.2).
\end{sesionbox}

\begin{sesionbox}{17}{Tarea}{mejora del libro y ensayo}
\campo{Inicio (5 min).} Cada equipo prioriza las críticas recibidas.

\campo{Desarrollo (30 min).} Se incorporan las críticas, se completan las partes pendientes y se cierra el registro de cambios. Cada miembro prepara la parte que va a exponer.

\campo{Ensayo (20 min).} Cada equipo ensaya ante otro, con tiempo controlado, y recibe una mejora concreta.

\campo{Agrupamiento:} equipos. \campo{Recursos:} libros, R2. \campo{Evaluación:} registro de cambios (9.2); observación (9.3). El libro se entrega al terminar la sesión.
\end{sesionbox}

\begin{sesionbox}{18}{Evaluación}{examen y retroalimentación}
\campo{Examen (45 min).} Prueba individual escrita (sección~\ref{sec:examen}).

\campo{Cierre (10 min).} Autocorrección guiada de un ejercicio; la devolución individual se entrega antes de S20.

\campo{Agrupamiento:} individual. \campo{Recursos:} examen, calculadora. \campo{Evaluación:} examen (1.1, 1.2, 2.1, 3.1, 5.1, 5.2, 7.2).
\end{sesionbox}

\begin{sesionbox}{19}{Cierre}{defendemos la propuesta (I)}
\campo{Apertura (3 min).} Orden de intervención y equipos asignados para coevaluar.

\campo{Defensas (32 min).} Cuatro equipos ante la clase y la dirección o un representante: 5 minutos de exposición, 2 de preguntas y 1 de cambio.

\campo{Coevaluación (15 min).} Cada equipo asignado completa la rúbrica del que ha valorado.

\campo{Cierre (5 min).} Primeras conclusiones sobre las propuestas.

\campo{Agrupamiento:} equipos y gran grupo. \campo{Recursos:} libros, proyector, rúbricas. \campo{Evaluación:} exposición (1.1, 2.2, 6.1, 6.2, 8.1); coevaluación (7.2, 8.1).
\end{sesionbox}

\begin{sesionbox}{20}{Cierre}{defendemos la propuesta (II) y cerramos}
\campo{Defensas (24 min).} Los tres equipos restantes, con el mismo formato.

\campo{Coevaluación (10 min).} Rúbrica de los equipos asignados y valoración dentro del propio equipo (9.1, 9.2, 9.3).

\campo{Metacognición (13 min).} Se retoman el mural de S1 y las estimaciones iniciales; cada estudiante completa la autoevaluación y cada equipo, la hoja de relevo del turno~4.

\campo{Cierre (8 min).} Videomensaje conjunto para el socio eTwinning.

\campo{Agrupamiento:} equipos y gran grupo. \campo{Recursos:} libros, proyector, rúbricas. \campo{Evaluación:} exposición (1.1, 2.2, 6.1, 6.2, 8.1); co/autoevaluación (7.2, 8.1, 9.1, 9.2, 9.3).
\end{sesionbox}
